\documentclass[11pt,a4paper]{article}

\usepackage[T1]{fontenc}
\usepackage[utf8]{inputenc}
\usepackage[polish]{babel}
\usepackage{lmodern}
\usepackage{graphicx}
\usepackage[table]{xcolor}
\usepackage{tabularx}
\usepackage{array}
\usepackage{geometry}
\usepackage{hyperref}
\usepackage{enumitem}
\usepackage{fancyhdr}
\usepackage{listings}
\usepackage{tcolorbox}

\geometry{left=2cm,right=2cm,top=1.7cm,bottom=1.8cm}
\setlength{\parindent}{0pt}
\setlength{\parskip}{5pt}

\definecolor{PSblue}{HTML}{1F5B99}
\definecolor{PSlight}{HTML}{EAF2F8}
\definecolor{PSgray}{HTML}{5F6B76}

\hypersetup{colorlinks=true,urlcolor=PSblue,linkcolor=PSblue}
\pagestyle{fancy}
\fancyhf{}
\fancyfoot[C]{\small\color{PSgray}Aplikacje webowe i sieciowe -- Politechnika Śląska}
\renewcommand{\headrulewidth}{0pt}

\newcommand{\pole}[1]{\textit{#1}}
\newcommand{\zadanie}{Laboratorium 3 -- Podstawy React + TypeScript + Vite}
\newcommand{\punkty}{10}
\newcommand{\skladsekcji}{Nazar Dovhopol}

\begin{document}

\begin{minipage}[c]{0.25\textwidth}
    \IfFileExists{LogoPSl.png}{\includegraphics[width=3.0cm]{LogoPSl.png}}{\fbox{\parbox[c][2.7cm][c]{3cm}{\centering Logo PŚ}}}
\end{minipage}%
\begin{minipage}[c]{0.75\textwidth}
    \raggedleft
    {\Large\bfseries\color{PSblue} APLIKACJE WEBOWE I SIECIOWE}\\[2mm]
    {\small\color{PSgray}Formatka dokumentacji zadania}
\end{minipage}

\vspace{2mm}
{\color{PSblue}\rule{\textwidth}{1.2pt}}
\vspace{4mm}

\begin{center}
{\LARGE\bfseries\color{PSblue} Dokumentacja wykonania zadania}\\[1mm]
{\small\color{PSgray}Szablon do zadań punktowanych realizowanych podczas laboratoriów}
\end{center}

\vspace{3mm}

\renewcommand{\arraystretch}{1.45}
\begin{tabularx}{\textwidth}{>{\bfseries\color{PSblue}\columncolor{PSlight}}p{4.6cm} X}
Przedmiot & Aplikacje webowe i sieciowe \\
Kierunek studiów & Informatyka \\
Forma i poziom studiów & studia stacjonarne I stopnia (inżynierskie), semestr 5 \\
Prowadzący & dr inż. Rafał Brociek \\
Skład sekcji / Imię i nazwisko & \skladsekcji \\
Zadanie & \zadanie \\
Liczba punktów & \punkty \\
\end{tabularx}

\section*{\color{PSblue}1. Zakres wykonanego zadania}
{\small\color{PSgray}\textit{Krótko opisz, co zostało wykonane w ramach zadania. Wskaż najważniejsze elementy rozwiązania i uzyskany rezultat.}}

\begin{tcolorbox}[colback=gray!2,colframe=gray!25,boxrule=0.5pt,arc=1mm]
W ramach laboratorium utworzyłem aplikację kliencką \texttt{Cars Client} w technologii React + TypeScript + Vite, wykorzystującą serwer \texttt{json-server} do obsługi operacji CRUD:
\begin{itemize}[leftmargin=*,itemsep=1pt,topsep=2pt]
    \item Skonfigurowałem projekt frontendowy \texttt{cars-client} w oparciu o Vite, TypeScript oraz pakiet \texttt{react-router}.
    \item Przygotowałem lokalny serwer danych \texttt{json-server} udostępniający bazę \texttt{db.json} z zasobem \texttt{/cars} na porcie 3000.
    \item Podzieliłem architekturę aplikacji na katalogi: \texttt{components/}, \texttt{layouts/}, \texttt{models/}, \texttt{pages/} oraz \texttt{services/}.
    \item Utworzyłem model danych \texttt{Car} oraz typ \texttt{NewCar} wykorzystujący \texttt{Omit<Car, 'id'>}.
    \item Zaimplementowałem dedykowany serwis \texttt{carService.ts} realizujący zapytania HTTP za pomocą wbudowanego \texttt{fetch} z obsługą \texttt{ensureSuccess} oraz sygnałów anulowania \texttt{AbortSignal}.
    \item Skonfigurowałem routing aplikacji w \texttt{App.tsx} oparty o \texttt{BrowserRouter}, wspólny layout \texttt{RootLayout} z aktywnym menu \texttt{NavLink} oraz widoki: \texttt{HomePage}, \texttt{CarsPage}, \texttt{CarDetailsPage}, \texttt{CreateCarPage}, \texttt{EditCarPage} i \texttt{NotFoundPage}.
    \item Zaimplementowałem formularze kontrolowane dla operacji dodawania (\texttt{POST}) oraz edycji (\texttt{PUT}) z wykorzystaniem nowoczesnych typów zdarzeń \texttt{SubmitEvent} i \texttt{ChangeEvent}.
    \item Zrealizowałem usuwanie samochodu z wymaganym potwierdzeniem użytkownika (\texttt{window.confirm}) oraz dynamiczną aktualizacją stanu.
    \item Wykonałem zadanie dodatkowe: dynamiczne filtrowanie listy samochodów po marce lub modelu po stronie klienta.
\end{itemize}
\end{tcolorbox}

\section*{\color{PSblue}2. Dokumentacja rozwiązania}
{\small\color{PSgray}\textit{W razie potrzeby zamieść kluczowe fragmenty kodu, krótkie objaśnienia, zrzuty ekranu lub inne elementy potwierdzające wykonanie zadania. Nie wklejaj całego kodu źródłowego.}}

\begin{tcolorbox}[colback=gray!2,colframe=gray!25,boxrule=0.5pt,arc=1mm]
\textbf{Część I: Architektura, komunikacja HTTP i routing}

\textbf{1. Model danych i lokalny serwer json-server} \\
W pliku \texttt{src/models/Car.ts} zdefiniowałem interfejs \texttt{Car} ze wszystkimi polami zgodnymi z modelem domenowym (marka, model, liczba drzwi, pojemność bagażnika, pojemność silnika, rodzaj paliwa, data produkcji, średnie spalanie, typ nadwozia) oraz typ pomocniczy \texttt{NewCar = Omit<Car, 'id'>}. Baza testowa została przygotowana w pliku \texttt{db.json}, a jej uruchomienie zautomatyzowane skryptem \texttt{npm run api}.

\textbf{2. Warstwa serwisowa carService.ts} \\
Komunikację z API odseparowałem od komponentów interfejsu. Funkcje \texttt{getCars}, \texttt{getCarById}, \texttt{createCar}, \texttt{updateCar} i \texttt{deleteCar} korzystają z funkcji \texttt{fetch}. Błędy HTTP weryfikowane są przez funkcję pomocniczą \texttt{ensureSuccess}:
\begin{lstlisting}[language=Java,basicstyle=\ttfamily\scriptsize,breaklines=true]
async function ensureSuccess(response: Response) {
  if (!response.ok) throw new Error(`Request failed. Status: ${response.status}`);
}
\end{lstlisting}
Zapytania pobierające dane przyjmują opcjonalny parametr \texttt{signal?: AbortSignal}, co pozwala na ich bezpieczne anulowanie w funkcji czyszczącej hooka \texttt{useEffect}.

\textbf{3. Drzewo tras i RootLayout} \\
W pliku \texttt{App.tsx} zdefiniowałem strukturę tras wewnątrz \texttt{RootLayout}, zawierającego stałe menu nawigacyjne z odnośnikami \texttt{NavLink} oraz kontener \texttt{<Outlet />}:
\begin{itemize}[leftmargin=*,itemsep=1pt,topsep=1pt]
    \item \texttt{/} -- strona główna (\texttt{HomePage}),
    \item \texttt{/cars} -- lista pojazdów wraz z wyszukiwarką (\texttt{CarsPage}),
    \item \texttt{/cars/new} -- formularz dodawania nowego auta (\texttt{CreateCarPage}),
    \item \texttt{/cars/:id} -- karta szczegółów pojazdu (\texttt{CarDetailsPage}),
    \item \texttt{/cars/:id/edit} -- formularz edycji istniejącego auta (\texttt{EditCarPage}),
    \item \texttt{*} -- strona błędu 404 (\texttt{NotFoundPage}).
\end{itemize}
\end{tcolorbox}

\newpage

\begin{tcolorbox}[colback=gray!2,colframe=gray!25,boxrule=0.5pt,arc=1mm]
\textbf{Część II: Formularze kontrolowane, operacje CRUD i zrzuty ekranu}

\textbf{4. Obsługa stanu, formularzy i usuwania} \\
Widok \texttt{CarsPage} pobiera dane w hooku \texttt{useEffect} i obsługuje trzy stany: ładowanie, błąd pobierania oraz pustą listę. Przed wykonaniem żądania \texttt{DELETE} wywoływane jest okno potwierdzenia \texttt{window.confirm}. Formularze w \texttt{CreateCarPage} i \texttt{EditCarPage} zostały zrealizowane jako komponenty kontrolowane ze stanem w \texttt{useState}. Edycja pobiera aktualne dane samochodu i wysyła pełny obiekt zaktualizowanego pojazdu żądaniem \texttt{PUT}.

\vspace{3mm}
\begin{center}
\begin{minipage}[t]{0.48\textwidth}
    \centering
    \IfFileExists{cars_list.png}{%
        \includegraphics[width=\linewidth,height=5.0cm,keepaspectratio]{cars_list.png}%
    }{%
        \fbox{\parbox[c][4.4cm][c]{0.9\linewidth}{\centering \small \textbf{Zrzut 1: Lista samochodów}\\[2mm] Zapisz plik jako \texttt{cars_list.png}.\\[1mm] (Widok \texttt{/cars} z filtrowaniem i kartami aut.)}}%
    }
    \\[1.5mm]
    {\small\textbf{Rysunek 1:} Widok listy samochodów z filtrowaniem.}
\end{minipage}
\hfill
\begin{minipage}[t]{0.48\textwidth}
    \centering
    \IfFileExists{car_details.png}{%
        \includegraphics[width=\linewidth,height=5.0cm,keepaspectratio]{car_details.png}%
    }{%
        \fbox{\parbox[c][4.4cm][c]{0.9\linewidth}{\centering \small \textbf{Zrzut 2: Szczegóły auta}\\[2mm] Zapisz plik jako \texttt{car_details.png}.\\[1mm] (Widok \texttt{/cars/:id} ze wszystkimi parametrami.)}}%
    }
    \\[1.5mm]
    {\small\textbf{Rysunek 2:} Widok szczegółów wybranego samochodu.}
\end{minipage}

\vspace{4mm}

\begin{minipage}[t]{0.48\textwidth}
    \centering
    \IfFileExists{car_form.png}{%
        \includegraphics[width=\linewidth,height=5.0cm,keepaspectratio]{car_form.png}%
    }{%
        \fbox{\parbox[c][4.4cm][c]{0.9\linewidth}{\centering \small \textbf{Zrzut 3: Formularz}\\[2mm] Zapisz plik jako \texttt{car_form.png}.\\[1mm] (Formularz dodawania \texttt{/cars/new} lub edycji \texttt{/cars/:id/edit}.)}}%
    }
    \\[1.5mm]
    {\small\textbf{Rysunek 3:} Kontrolowany formularz dodawania/edycji.}
\end{minipage}
\hfill
\begin{minipage}[t]{0.48\textwidth}
    \centering
    \IfFileExists{delete_confirm.png}{%
        \includegraphics[width=\linewidth,height=5.0cm,keepaspectratio]{delete_confirm.png}%
    }{%
        \fbox{\parbox[c][4.4cm][c]{0.9\linewidth}{\centering \small \textbf{Zrzut 4: Potwierdzenie usuwania}\\[2mm] Zapisz plik jako \texttt{delete_confirm.png}.\\[1mm] (Okno \texttt{window.confirm} przed usunięciem lub działający \texttt{json-server}.)}}%
    }
    \\[1.5mm]
    {\small\textbf{Rysunek 4:} Okno dialogowe potwierdzenia usunięcia auta.}
\end{minipage}
\end{center}
\end{tcolorbox}

\section*{\color{PSblue}3. Repozytorium GitHub}
{\small\color{PSgray}\textit{Podaj bezpośredni link do repozytorium zawierającego rozwiązanie zadania.}}

\begin{tabularx}{\textwidth}{>{\bfseries\color{PSblue}\columncolor{PSlight}}p{2.3cm} X}
Link: & \url{https://github.com/Waterlo-o/UniversityCars} \\
\end{tabularx}

\vfill
\begin{flushright}
{\small\color{PSgray}Data przygotowania: \today}
\end{flushright}

\end{document}
